\begin{claim}{D4.3}{Proved}{The exceptional nonintegral twists cannot be omitted}
On a regular marked fibre, $M_z$ is trivial if and only if there are
integers $n,p,q,r$ such that
\[
 q\tau(t)+r\mu(t)=p,\qquad
 z=n-6q\mu(t)-r\beta_c(t).                             \tag{D.15}
\]
For $z\notin\Z$ these resonances occur in a countable set of parameter
values $z$; every resonant $z$ is nonreal. On the finite root covers
the exceptional lattices are
\[
 \Lambda_1=\Z+(3\beta_c(z_1)+4-2\rho)\Z,\qquad
 \Lambda_2=\Z+(2\beta_c(z_2)+3-3i)\Z.                  \tag{D.16}
\]
Take the union of $\Z$, these two lattices, and all values from (D.15)
over regular fibres; this is the countable exceptional set used in
D4.2. The set is larger than $\Z$: the regular value $t=32/5$
supplies nonintegral twists with $h^1(X,T_X\otimes M_z)\ge2$.
\end{claim}
\begin{proof}
The exponential sequence on a torus identifies trivial line bundles
with integral constant cocycles modulo the two holomorphic period
rows. Write
$z\gamma=(n,p,q,r)+a(6\mu,\tau,1,0)+b(\beta_c,\mu,0,1)$.
The third and fourth coordinates give $a=-q,b=-r$; the first two
give exactly (D.15). If $r=0$, the positivity of $\operatorname{Im}\tau$
forces $q=p=0$, leaving $z=n$. Otherwise the imaginary parts of
(D.15) give
\[
 \operatorname{Im}\mu=-\frac qr\operatorname{Im}\tau,
 \qquad\operatorname{Im}z=-rD\ne0.                    \tag{D.17}
\]
For fixed $(p,q,r)$ with $(q,r)\ne(0,0)$, the first equation of
(D.15) has discrete zeros, since $1,\tau,\mu$ are linearly
independent. A discrete subset of this second-countable Riemann
surface is countable. There are countably many integer tuples and
markings, so the union of the resulting $z$ is countable.

At $z_1$, insert $\tau=\rho$, $\mu=(2-\rho)/3$ in the first
equation. Its solutions have $(q,r,p)=(k,3k,2k)$.
At $z_2$ they have $(q,r,p)=(k,2k,k)$.
The second equation gives (D.16). These are discrete lattices because
their nonintegral generators have imaginary parts $3D$ and $2D$.

For the explicit regular resonance, use the section
$P=(0,\sqrt2t^2(t-1))$ of the stated elliptic surface. The
chord-and-tangent law gives
\[
 x(2P)=9t^2/8,\qquad
 y(2P)=t^2(32-5t)/(16\sqrt2).
\]
At $t=32/5$, the discriminant
$1728t^8(t-1)^3$ is nonzero, $P$ is not two-torsion, and $2P$ is a
nonzero two-torsion point. Thus $P$ has exact order four. The period
identification in Part A supplies a nonzero relation
$q\tau+4\mu=p$ with integers $p,q$. Equation (D.15), with $r=4$,
defines an actual twist for which the regular torus restriction is
trivial; (D.17) shows that it is nonintegral. The period derivative
there has nonzero determinant, so the tangent restriction has
$h^0(T_X|_F)=2$. The degree-zero base-change exact sequence gives
\[
 H^0(F,(T_X\otimes M_z)|_F)
   \simeq (R^1f_*(T_X\otimes M_z))_t[t],
\]
because $f_*(T_X\otimes M_z)=0$ for a twist nontrivial on the general
fibre. Its right side has dimension two. The degree-one direct image
is a finite-support sheaf, so its global sections have dimension at
least two, proving the bound on $h^1$. The exact elliptic calculation
is checked by both duplication and substitution in
\texttt{d\_exact.py}; its output records the nonzero rational
discriminant. The choice of integers $p,q$ depends on the marked
elliptic logarithm, and no unverified numerical marking is substituted.
\end{proof}
